
\documentclass{article}
\usepackage{colortbl}
\usepackage{makecell}
\usepackage{multirow}
\usepackage{supertabular}

\begin{document}

\newcounter{utterance}

\centering \large Interaction Transcript for game `clean\_up', experiment `2\_hard\_3obj\_nl', episode 3 with deepseek{-}v4{-}pro.
\vspace{24pt}

{ \footnotesize  \setcounter{utterance}{1}
\setlength{\tabcolsep}{0pt}
\begin{supertabular}{c@{$\;$}|p{.15\linewidth}@{}p{.15\linewidth}p{.15\linewidth}p{.15\linewidth}p{.15\linewidth}p{.15\linewidth}}
    \# & $\;$A & \multicolumn{4}{c}{Game Master} & $\;\:$B\\
    \hline

    \theutterance \stepcounter{utterance}  
    & & \multicolumn{4}{p{0.6\linewidth}}{
        \cellcolor[rgb]{0.9,0.9,0.9}{
            \makecell[{{p{\linewidth}}}]{
                \texttt{\tiny{[P1$\langle$GM]}}
                \texttt{Ik ben uw spelleider, en u speelt een samenwerkingsspel met het volgende raster als speelbord:} \\
\\ 
\texttt{```} \\
\texttt{ 1234567} \\
\texttt{╔═╤═╤═══╗} \\
\texttt{║L│◌│◌◌◌║ 1} \\
\texttt{║C├─┤◌◌◌║ 2} \\
\texttt{╟─┤◌├─┬─╢ 3} \\
\texttt{║◌│◌│◌│◌║ 4} \\
\texttt{╟─┼─┘◌│◌║ 5} \\
\texttt{║◌│◌◌┌┴─╢ 6} \\
\texttt{║◌│◌◌│◌P║ 7} \\
\texttt{╚═╧══╧══╝} \\
\\ 
\texttt{```} \\
\\ 
\texttt{* De bovenrand toont x{-}coördinaten die naar rechts oplopen, en de rechterrand y{-}coördinaten die naar beneden oplopen.} \\
\texttt{* De volgende objecten zijn willekeurig op uw raster geplaatst: 'C', 'L', 'P'.} \\
\\ 
\\ 
\texttt{De andere speler ziet een variatie van het speelbord, waarbij de objecten op verschillende willekeurige locaties zijn geplaatst. U kunt het bord van de andere speler niet zien en zij kunnen het uwe niet zien.} \\
\\ 
\\ 
\texttt{**Doel:**} \\
\\ 
\texttt{Beide spelers moeten de objecten op hun respectievelijke achtergrond zo verplaatsen dat identieke objecten op dezelfde coördinaten terechtkomen. U moet communiceren met de andere speler om een gemeenschappelijke doelconfiguratie overeen te komen.} \\
\\ 
\\ 
\texttt{**Regels:**} \\
\\ 
\texttt{* In elke beurt kunt u precies één van de volgende twee commando's sturen:} \\
\texttt{ 1. `ZEG: <BERICHT>`: om een bericht (alles tot de volgende regeleinde) naar de andere speler te sturen. Ik zal het doorsturen naar uw partner.} \\
\texttt{2. `BEWEEG: <OBJECT>, (<X>, <Y>)`: om een object naar een nieuwe positie te verplaatsen, waarbij `<X>` de kolom is en `<Y>` de rij. Ik zal u laten weten of uw zet succesvol was of niet.} \\
\texttt{* Als u zich niet aan het formaat houdt, of meerdere commando's tegelijk verstuurt, moet ik u straffen.} \\
\texttt{* Als beide spelers meer dan 8 strafpunten verzamelen, verliezen jullie allebei het spel.} \\
\texttt{* Het is van vitaal belang dat u met de andere speler communiceert over uw doelstatus! De *enige* manier waarop u uw strategie aan de andere speler kunt overbrengen is door gebruik te maken van het `ZEG: <BERICHT>` commando!} \\
\\ 
\\ 
\texttt{**Objecten verplaatsen**} \\
\\ 
\texttt{* U kunt alleen objecten verplaatsen naar cellen binnen de grenzen van het raster. De doelcel moet leeg zijn, d.w.z. deze mag alleen het symbool '◌' bevatten.} \\
\texttt{* Als u probeert een object te verplaatsen naar een plek die niet leeg is, of als u probeert het buiten het raster te verplaatsen, moet ik u straffen. U krijgt een nieuwe poging.} \\
\texttt{* Voordat u een zet doet, moet u controleren of de doelplek leeg is, en geen letter, kader of lijn bevat!} \\
\\ 
\\ 
\texttt{**Einde van het spel**} \\
\\ 
\texttt{Als u denkt dat u het doel bereikt hebt om alle objecten uit te lijnen, kunt u de andere speler vragen het spel te beëindigen door `ZEG: klaar?` te sturen. Als de andere speler u vraagt om het spel te beëindigen en u antwoordt met `ZEG: klaar!`, zal het spel eindigen.} \\
\\ 
\texttt{Beide spelers winnen als het spel eindigt binnen 12 rondes, waarbij één ronde gedefinieerd wordt als twee spelers die elk een geldig commando sturen.} \\
\\ 
\texttt{**Score:**} \\
\\ 
\texttt{Hoe dichter de identieke objecten bij elkaar liggen op beide speelborden, hoe meer punten u krijgt. Strafpunten verminderen uw punten. Kunt u het record verslaan?} \\
\\ 
\\ 
\texttt{Stuur een bericht naar de andere speler om het spel te starten!} \\
            }
        }
    }
    & \\ \\

    \theutterance \stepcounter{utterance}  
    & \multicolumn{4}{p{0.6\linewidth}}{
        \cellcolor[rgb]{0.9,0.9,0.9}{
            \makecell[{{p{\linewidth}}}]{
                \texttt{\tiny{[P1$\rangle$GM]}}
                \texttt{ZEG: Hallo! Laten we onze posities delen om een gemeenschappelijke doelconfiguratie af te spreken. Mijn objecten staan op: C=(2,2), L=(1,1), P=(7,7). Wat zijn jouw coördinaten?} \\
            }
        }
    }
    & & \\ \\

    \theutterance \stepcounter{utterance}  
    & & & \multicolumn{2}{p{0.3\linewidth}}{
        \cellcolor[rgb]{0.9,0.9,0.9}{
            \makecell[{{p{\linewidth}}}]{
                \texttt{\tiny{[GM$|$GM]}}
                \texttt{ZEG: Hallo! Laten we onze posities delen om een gemeenschappelijke doelconfiguratie af te spreken. Mijn objecten staan op: C=(2,2), L=(1,1), P=(7,7). Wat zijn jouw coördinaten?} \\
            }
        }
    }
    & & \\ \\

    \theutterance \stepcounter{utterance}  
    & & \multicolumn{4}{p{0.6\linewidth}}{
        \cellcolor[rgb]{0.9,0.9,0.9}{
            \makecell[{{p{\linewidth}}}]{
                \texttt{\tiny{[GM$\rangle$P2]}}
                \texttt{Ik ben uw spelleider, en u speelt een samenwerkingsspel met het volgende raster als speelbord:} \\
\\ 
\texttt{```} \\
\texttt{ 1234567} \\
\texttt{╔═╤═╤═══╗} \\
\texttt{║◌│◌│◌◌◌║ 1} \\
\texttt{║◌├─┤◌L◌║ 2} \\
\texttt{╟─┤P├─┬─╢ 3} \\
\texttt{║◌│◌│◌│◌║ 4} \\
\texttt{╟─┼─┘C│◌║ 5} \\
\texttt{║◌│◌◌┌┴─╢ 6} \\
\texttt{║◌│◌◌│◌◌║ 7} \\
\texttt{╚═╧══╧══╝} \\
\\ 
\texttt{```} \\
\\ 
\texttt{* De bovenrand toont x{-}coördinaten die naar rechts oplopen, en de rechterrand y{-}coördinaten die naar beneden oplopen.} \\
\texttt{* De volgende objecten zijn willekeurig op uw raster geplaatst: 'C', 'L', 'P'.} \\
\\ 
\\ 
\texttt{De andere speler ziet een variatie van het speelbord, waarbij de objecten op verschillende willekeurige locaties zijn geplaatst. U kunt het bord van de andere speler niet zien en zij kunnen het uwe niet zien.} \\
\\ 
\\ 
\texttt{**Doel:**} \\
\\ 
\texttt{Beide spelers moeten de objecten op hun respectievelijke achtergrond zo verplaatsen dat identieke objecten op dezelfde coördinaten terechtkomen. U moet communiceren met de andere speler om een gemeenschappelijke doelconfiguratie overeen te komen.} \\
\\ 
\\ 
\texttt{**Regels:**} \\
\\ 
\texttt{* In elke beurt kunt u precies één van de volgende twee commando's sturen:} \\
\texttt{ 1. `ZEG: <BERICHT>`: om een bericht (alles tot de volgende regeleinde) naar de andere speler te sturen. Ik zal het doorsturen naar uw partner.} \\
\texttt{2. `BEWEEG: <OBJECT>, (<X>, <Y>)`: om een object naar een nieuwe positie te verplaatsen, waarbij `<X>` de kolom is en `<Y>` de rij. Ik zal u laten weten of uw zet succesvol was of niet.} \\
\texttt{* Als u zich niet aan het formaat houdt, of meerdere commando's tegelijk verstuurt, moet ik u straffen.} \\
\texttt{* Als beide spelers meer dan 8 strafpunten verzamelen, verliezen jullie allebei het spel.} \\
\texttt{* Het is van vitaal belang dat u met de andere speler communiceert over uw doelstatus! De *enige* manier waarop u uw strategie aan de andere speler kunt overbrengen is door gebruik te maken van het `ZEG: <BERICHT>` commando!} \\
\\ 
\\ 
\texttt{**Objecten verplaatsen**} \\
\\ 
\texttt{* U kunt alleen objecten verplaatsen naar cellen binnen de grenzen van het raster. De doelcel moet leeg zijn, d.w.z. deze mag alleen het symbool '◌' bevatten.} \\
\texttt{* Als u probeert een object te verplaatsen naar een plek die niet leeg is, of als u probeert het buiten het raster te verplaatsen, moet ik u straffen. U krijgt een nieuwe poging.} \\
\texttt{* Voordat u een zet doet, moet u controleren of de doelplek leeg is, en geen letter, kader of lijn bevat!} \\
\\ 
\\ 
\texttt{**Einde van het spel**} \\
\\ 
\texttt{Als u denkt dat u het doel bereikt hebt om alle objecten uit te lijnen, kunt u de andere speler vragen het spel te beëindigen door `ZEG: klaar?` te sturen. Als de andere speler u vraagt om het spel te beëindigen en u antwoordt met `ZEG: klaar!`, zal het spel eindigen.} \\
\\ 
\texttt{Beide spelers winnen als het spel eindigt binnen 12 rondes, waarbij één ronde gedefinieerd wordt als twee spelers die elk een geldig commando sturen.} \\
\\ 
\texttt{**Score:**} \\
\\ 
\texttt{Hoe dichter de identieke objecten bij elkaar liggen op beide speelborden, hoe meer punten u krijgt. Strafpunten verminderen uw punten. Kunt u het record verslaan?} \\
\\ 
\texttt{De andere speler begon het spel door dit bericht te sturen:} \\
\\ 
\texttt{"Hallo! Laten we onze posities delen om een gemeenschappelijke doelconfiguratie af te spreken. Mijn objecten staan op: C=(2,2), L=(1,1), P=(7,7). Wat zijn jouw coördinaten?"} \\
\\ 
\texttt{Wat is uw eerste opdracht?} \\
            }
        }
    }
    & \\ \\

    \theutterance \stepcounter{utterance}  
    & & & \multicolumn{4}{p{0.6\linewidth}}{
        \cellcolor[rgb]{0.9,0.9,0.9}{
            \makecell[{{p{\linewidth}}}]{
                \texttt{\tiny{[GM$\langle$P2]}}
                \texttt{Mijn objecten staan op de volgende coördinaten:} \\
\\ 
\texttt{{-} C = (4,5)} \\
\texttt{{-} L = (5,2)} \\
\texttt{{-} P = (3,3)} \\
\\ 
\texttt{Dus laten we afspreken waar we ze naartoe verplaatsen!} \\
\\ 
\texttt{`ZEG: Hallo! Mijn objecten staan op: C=(4,5), L=(5,2), P=(3,3). Laten we een gemeenschappelijke doelconfiguratie afspreken. Welke posities stel je voor?`} \\
            }
        }
    }
     \\ \\

    \theutterance \stepcounter{utterance}  
    & & & \multicolumn{2}{p{0.3\linewidth}}{
        \cellcolor[rgb]{0.9,0.9,0.9}{
            \makecell[{{p{\linewidth}}}]{
                \texttt{\tiny{[GM$|$GM]}}
                \texttt{Mijn objecten staan op de volgende coördinaten:} \\
\\ 
\texttt{{-} C = (4,5)} \\
\texttt{{-} L = (5,2)} \\
\texttt{{-} P = (3,3)} \\
\\ 
\texttt{Dus laten we afspreken waar we ze naartoe verplaatsen!} \\
\\ 
\texttt{`ZEG: Hallo! Mijn objecten staan op: C=(4,5), L=(5,2), P=(3,3). Laten we een gemeenschappelijke doelconfiguratie afspreken. Welke posities stel je voor?`} \\
            }
        }
    }
    & & \\ \\

    \theutterance \stepcounter{utterance}  
    & & & \multicolumn{2}{p{0.3\linewidth}}{
        \cellcolor[rgb]{0.9,0.9,0.9}{
            \makecell[{{p{\linewidth}}}]{
                \texttt{\tiny{[GM$|$GM]}}
                \texttt{Invalid format: head is not empty: 'Mijn objecten staan op de volgende coördinaten:} \\
\\ 
\texttt{{-} C = (4,5)} \\
\texttt{{-} L = (5,2)} \\
\texttt{{-} P = (3,3)} \\
\\ 
\texttt{Dus laten we afspreken waar we ze naartoe verplaatsen!} \\
\\ 
\texttt{'} \\
            }
        }
    }
    & & \\ \\

    \theutterance \stepcounter{utterance}  
    & & \multicolumn{4}{p{0.6\linewidth}}{
        \cellcolor[rgb]{0.9,0.9,0.9}{
            \makecell[{{p{\linewidth}}}]{
                \texttt{\tiny{[GM$\rangle$P2]}}
                \texttt{Straf: Je bericht mag niets vóór het commando bevatten!} \\
\\ 
\texttt{Zorg ervoor dat je antwoord alleen `ZEG: <BERICHT>` of `BEWEEG: <OBJECT>, (<X>, <Y>)` bevat, en niets anders!} \\
\\ 
\texttt{Jullie hebben gezamenlijk 1 van 8 straffen verzameld.} \\
\texttt{Probeer het opnieuw!} \\
            }
        }
    }
    & \\ \\

    \theutterance \stepcounter{utterance}  
    & & & \multicolumn{4}{p{0.6\linewidth}}{
        \cellcolor[rgb]{0.9,0.9,0.9}{
            \makecell[{{p{\linewidth}}}]{
                \texttt{\tiny{[GM$\langle$P2]}}
                \texttt{`ZEG: Hallo! Mijn objecten staan op: C=(4,5), L=(5,2), P=(3,3). Laten we een gemeenschappelijke doelconfiguratie afspreken. Welke posities stel je voor?`} \\
            }
        }
    }
     \\ \\

    \theutterance \stepcounter{utterance}  
    & & & \multicolumn{2}{p{0.3\linewidth}}{
        \cellcolor[rgb]{0.9,0.9,0.9}{
            \makecell[{{p{\linewidth}}}]{
                \texttt{\tiny{[GM$|$GM]}}
                \texttt{`ZEG: Hallo! Mijn objecten staan op: C=(4,5), L=(5,2), P=(3,3). Laten we een gemeenschappelijke doelconfiguratie afspreken. Welke posities stel je voor?`} \\
            }
        }
    }
    & & \\ \\

    \theutterance \stepcounter{utterance}  
    & & \multicolumn{4}{p{0.6\linewidth}}{
        \cellcolor[rgb]{0.9,0.9,0.9}{
            \makecell[{{p{\linewidth}}}]{
                \texttt{\tiny{[P1$\langle$GM]}}
                \texttt{Je bericht is doorgestuurd naar de andere speler.} \\
\\ 
\\ 
\texttt{Je speelt momenteel ronde 1 van maximaal 12.} \\
\texttt{Jullie hebben gezamenlijk 1 van 8 straffen verzameld.} \\
\texttt{De andere speler stuurde dit bericht:} \\
\\ 
\\ 
\texttt{"Hallo! Mijn objecten staan op: C=(4,5), L=(5,2), P=(3,3). Laten we een gemeenschappelijke doelconfiguratie afspreken. Welke posities stel je voor?"} \\
\\ 
\\ 
\texttt{Wat is je volgende opdracht?} \\
            }
        }
    }
    & \\ \\

    \theutterance \stepcounter{utterance}  
    & \multicolumn{4}{p{0.6\linewidth}}{
        \cellcolor[rgb]{0.9,0.9,0.9}{
            \makecell[{{p{\linewidth}}}]{
                \texttt{\tiny{[P1$\rangle$GM]}}
                \texttt{ZEG: Laten we een centrale locatie kiezen die voor ons beiden bereikbaar is. Wat dacht je van: C=(4,4), L=(3,3), P=(5,5)? Dan moeten we allebei wat bewegen, maar het is symmetrisch. Of heb jij een ander voorstel?} \\
            }
        }
    }
    & & \\ \\

    \theutterance \stepcounter{utterance}  
    & & & \multicolumn{2}{p{0.3\linewidth}}{
        \cellcolor[rgb]{0.9,0.9,0.9}{
            \makecell[{{p{\linewidth}}}]{
                \texttt{\tiny{[GM$|$GM]}}
                \texttt{ZEG: Laten we een centrale locatie kiezen die voor ons beiden bereikbaar is. Wat dacht je van: C=(4,4), L=(3,3), P=(5,5)? Dan moeten we allebei wat bewegen, maar het is symmetrisch. Of heb jij een ander voorstel?} \\
            }
        }
    }
    & & \\ \\

    \theutterance \stepcounter{utterance}  
    & & \multicolumn{4}{p{0.6\linewidth}}{
        \cellcolor[rgb]{0.9,0.9,0.9}{
            \makecell[{{p{\linewidth}}}]{
                \texttt{\tiny{[GM$\rangle$P2]}}
                \texttt{Je bericht is doorgestuurd naar de andere speler.} \\
\\ 
\\ 
\texttt{Je speelt momenteel ronde 2 van maximaal 12.} \\
\texttt{Jullie hebben gezamenlijk 1 van 8 straffen verzameld.} \\
\texttt{De andere speler stuurde dit bericht:} \\
\\ 
\\ 
\texttt{"Laten we een centrale locatie kiezen die voor ons beiden bereikbaar is. Wat dacht je van: C=(4,4), L=(3,3), P=(5,5)? Dan moeten we allebei wat bewegen, maar het is symmetrisch. Of heb jij een ander voorstel?"} \\
\\ 
\\ 
\texttt{Wat is je volgende opdracht?} \\
            }
        }
    }
    & \\ \\

    \theutterance \stepcounter{utterance}  
    & & & \multicolumn{4}{p{0.6\linewidth}}{
        \cellcolor[rgb]{0.9,0.9,0.9}{
            \makecell[{{p{\linewidth}}}]{
                \texttt{\tiny{[GM$\langle$P2]}}
                \texttt{`ZEG: C=(4,4) werkt voor mij, maar L=(3,3) is bezet door mijn P. Mijn P staat op (3,3). Ik stel voor: C=(4,4), L=(2,2), P=(6,6). Wat denk jij?`} \\
            }
        }
    }
     \\ \\

    \theutterance \stepcounter{utterance}  
    & & & \multicolumn{2}{p{0.3\linewidth}}{
        \cellcolor[rgb]{0.9,0.9,0.9}{
            \makecell[{{p{\linewidth}}}]{
                \texttt{\tiny{[GM$|$GM]}}
                \texttt{`ZEG: C=(4,4) werkt voor mij, maar L=(3,3) is bezet door mijn P. Mijn P staat op (3,3). Ik stel voor: C=(4,4), L=(2,2), P=(6,6). Wat denk jij?`} \\
            }
        }
    }
    & & \\ \\

    \theutterance \stepcounter{utterance}  
    & & \multicolumn{4}{p{0.6\linewidth}}{
        \cellcolor[rgb]{0.9,0.9,0.9}{
            \makecell[{{p{\linewidth}}}]{
                \texttt{\tiny{[P1$\langle$GM]}}
                \texttt{Je bericht is doorgestuurd naar de andere speler.} \\
\\ 
\\ 
\texttt{Je speelt momenteel ronde 2 van maximaal 12.} \\
\texttt{Jullie hebben gezamenlijk 1 van 8 straffen verzameld.} \\
\texttt{De andere speler stuurde dit bericht:} \\
\\ 
\\ 
\texttt{"C=(4,4) werkt voor mij, maar L=(3,3) is bezet door mijn P. Mijn P staat op (3,3). Ik stel voor: C=(4,4), L=(2,2), P=(6,6). Wat denk jij?"} \\
\\ 
\\ 
\texttt{Wat is je volgende opdracht?} \\
            }
        }
    }
    & \\ \\

    \theutterance \stepcounter{utterance}  
    & \multicolumn{4}{p{0.6\linewidth}}{
        \cellcolor[rgb]{0.9,0.9,0.9}{
            \makecell[{{p{\linewidth}}}]{
                \texttt{\tiny{[P1$\rangle$GM]}}
                \texttt{ZEG: Klinkt goed! L=(2,2) is voor mij vrij want daar staat mijn C nu, maar die gaat naar (4,4). P=(6,6) is ook vrij bij mij. Laten we afspreken: C=(4,4), L=(2,2), P=(6,6). Ik begin met verplaatsen. Ben je akkoord?} \\
            }
        }
    }
    & & \\ \\

    \theutterance \stepcounter{utterance}  
    & & & \multicolumn{2}{p{0.3\linewidth}}{
        \cellcolor[rgb]{0.9,0.9,0.9}{
            \makecell[{{p{\linewidth}}}]{
                \texttt{\tiny{[GM$|$GM]}}
                \texttt{ZEG: Klinkt goed! L=(2,2) is voor mij vrij want daar staat mijn C nu, maar die gaat naar (4,4). P=(6,6) is ook vrij bij mij. Laten we afspreken: C=(4,4), L=(2,2), P=(6,6). Ik begin met verplaatsen. Ben je akkoord?} \\
            }
        }
    }
    & & \\ \\

    \theutterance \stepcounter{utterance}  
    & & \multicolumn{4}{p{0.6\linewidth}}{
        \cellcolor[rgb]{0.9,0.9,0.9}{
            \makecell[{{p{\linewidth}}}]{
                \texttt{\tiny{[GM$\rangle$P2]}}
                \texttt{Je bericht is doorgestuurd naar de andere speler.} \\
\\ 
\\ 
\texttt{Je speelt momenteel ronde 3 van maximaal 12.} \\
\texttt{Jullie hebben gezamenlijk 1 van 8 straffen verzameld.} \\
\texttt{De andere speler stuurde dit bericht:} \\
\\ 
\\ 
\texttt{"Klinkt goed! L=(2,2) is voor mij vrij want daar staat mijn C nu, maar die gaat naar (4,4). P=(6,6) is ook vrij bij mij. Laten we afspreken: C=(4,4), L=(2,2), P=(6,6). Ik begin met verplaatsen. Ben je akkoord?"} \\
\\ 
\\ 
\texttt{Wat is je volgende opdracht?} \\
            }
        }
    }
    & \\ \\

    \theutterance \stepcounter{utterance}  
    & & & \multicolumn{4}{p{0.6\linewidth}}{
        \cellcolor[rgb]{0.9,0.9,0.9}{
            \makecell[{{p{\linewidth}}}]{
                \texttt{\tiny{[GM$\langle$P2]}}
                \texttt{`ZEG: Akkoord! Doelconfiguratie: C=(4,4), L=(2,2), P=(6,6). Ik begin ook met verplaatsen.`} \\
            }
        }
    }
     \\ \\

    \theutterance \stepcounter{utterance}  
    & & & \multicolumn{2}{p{0.3\linewidth}}{
        \cellcolor[rgb]{0.9,0.9,0.9}{
            \makecell[{{p{\linewidth}}}]{
                \texttt{\tiny{[GM$|$GM]}}
                \texttt{`ZEG: Akkoord! Doelconfiguratie: C=(4,4), L=(2,2), P=(6,6). Ik begin ook met verplaatsen.`} \\
            }
        }
    }
    & & \\ \\

    \theutterance \stepcounter{utterance}  
    & & \multicolumn{4}{p{0.6\linewidth}}{
        \cellcolor[rgb]{0.9,0.9,0.9}{
            \makecell[{{p{\linewidth}}}]{
                \texttt{\tiny{[P1$\langle$GM]}}
                \texttt{Je bericht is doorgestuurd naar de andere speler.} \\
\\ 
\\ 
\texttt{Je speelt momenteel ronde 3 van maximaal 12.} \\
\texttt{Jullie hebben gezamenlijk 1 van 8 straffen verzameld.} \\
\texttt{De andere speler stuurde dit bericht:} \\
\\ 
\\ 
\texttt{"Akkoord! Doelconfiguratie: C=(4,4), L=(2,2), P=(6,6). Ik begin ook met verplaatsen."} \\
\\ 
\\ 
\texttt{Wat is je volgende opdracht?} \\
            }
        }
    }
    & \\ \\

    \theutterance \stepcounter{utterance}  
    & \multicolumn{4}{p{0.6\linewidth}}{
        \cellcolor[rgb]{0.9,0.9,0.9}{
            \makecell[{{p{\linewidth}}}]{
                \texttt{\tiny{[P1$\rangle$GM]}}
                \texttt{BEWEEG: C, (4,4)} \\
            }
        }
    }
    & & \\ \\

    \theutterance \stepcounter{utterance}  
    & & & \multicolumn{2}{p{0.3\linewidth}}{
        \cellcolor[rgb]{0.9,0.9,0.9}{
            \makecell[{{p{\linewidth}}}]{
                \texttt{\tiny{[GM$|$GM]}}
                \texttt{BEWEEG: C, (4,4)} \\
            }
        }
    }
    & & \\ \\

    \theutterance \stepcounter{utterance}  
    & & & \multicolumn{2}{p{0.3\linewidth}}{
        \cellcolor[rgb]{0.9,0.9,0.9}{
            \makecell[{{p{\linewidth}}}]{
                \texttt{\tiny{[GM$|$GM]}}
                \texttt{Straf: (4,4) is niet leeg, maar bevat '│'.} \\
\\ 
\texttt{Jullie hebben gezamenlijk 2 van 8 straffen verzameld.} \\
\texttt{Probeer het opnieuw!} \\
            }
        }
    }
    & & \\ \\

    \theutterance \stepcounter{utterance}  
    & & \multicolumn{4}{p{0.6\linewidth}}{
        \cellcolor[rgb]{0.9,0.9,0.9}{
            \makecell[{{p{\linewidth}}}]{
                \texttt{\tiny{[P1$\langle$GM]}}
                \texttt{Straf: (4,4) is niet leeg, maar bevat '│'.} \\
\\ 
\texttt{Jullie hebben gezamenlijk 2 van 8 straffen verzameld.} \\
\texttt{Probeer het opnieuw!} \\
            }
        }
    }
    & \\ \\

    \theutterance \stepcounter{utterance}  
    & \multicolumn{4}{p{0.6\linewidth}}{
        \cellcolor[rgb]{0.9,0.9,0.9}{
            \makecell[{{p{\linewidth}}}]{
                \texttt{\tiny{[P1$\rangle$GM]}}
                \texttt{BEWEEG: C, (2,4)} \\
            }
        }
    }
    & & \\ \\

    \theutterance \stepcounter{utterance}  
    & & & \multicolumn{2}{p{0.3\linewidth}}{
        \cellcolor[rgb]{0.9,0.9,0.9}{
            \makecell[{{p{\linewidth}}}]{
                \texttt{\tiny{[GM$|$GM]}}
                \texttt{BEWEEG: C, (2,4)} \\
            }
        }
    }
    & & \\ \\

    \theutterance \stepcounter{utterance}  
    & & & \multicolumn{2}{p{0.3\linewidth}}{
        \cellcolor[rgb]{0.9,0.9,0.9}{
            \makecell[{{p{\linewidth}}}]{
                \texttt{\tiny{[GM$|$GM]}}
                \texttt{Straf: (2,4) is niet leeg, maar bevat '│'.} \\
\\ 
\texttt{Jullie hebben gezamenlijk 3 van 8 straffen verzameld.} \\
\texttt{Probeer het opnieuw!} \\
            }
        }
    }
    & & \\ \\

    \theutterance \stepcounter{utterance}  
    & & \multicolumn{4}{p{0.6\linewidth}}{
        \cellcolor[rgb]{0.9,0.9,0.9}{
            \makecell[{{p{\linewidth}}}]{
                \texttt{\tiny{[P1$\langle$GM]}}
                \texttt{Straf: (2,4) is niet leeg, maar bevat '│'.} \\
\\ 
\texttt{Jullie hebben gezamenlijk 3 van 8 straffen verzameld.} \\
\texttt{Probeer het opnieuw!} \\
            }
        }
    }
    & \\ \\

    \theutterance \stepcounter{utterance}  
    & \multicolumn{4}{p{0.6\linewidth}}{
        \cellcolor[rgb]{0.9,0.9,0.9}{
            \makecell[{{p{\linewidth}}}]{
                \texttt{\tiny{[P1$\rangle$GM]}}
                \texttt{BEWEEG: C, (1,2)} \\
            }
        }
    }
    & & \\ \\

    \theutterance \stepcounter{utterance}  
    & & & \multicolumn{2}{p{0.3\linewidth}}{
        \cellcolor[rgb]{0.9,0.9,0.9}{
            \makecell[{{p{\linewidth}}}]{
                \texttt{\tiny{[GM$|$GM]}}
                \texttt{BEWEEG: C, (1,2)} \\
            }
        }
    }
    & & \\ \\

    \theutterance \stepcounter{utterance}  
    & & & \multicolumn{2}{p{0.3\linewidth}}{
        \cellcolor[rgb]{0.9,0.9,0.9}{
            \makecell[{{p{\linewidth}}}]{
                \texttt{\tiny{[GM$|$GM]}}
                \texttt{Straf: (1,2) is niet leeg, maar bevat 'C'.} \\
\\ 
\texttt{Jullie hebben gezamenlijk 4 van 8 straffen verzameld.} \\
\texttt{Probeer het opnieuw!} \\
            }
        }
    }
    & & \\ \\

    \theutterance \stepcounter{utterance}  
    & & \multicolumn{4}{p{0.6\linewidth}}{
        \cellcolor[rgb]{0.9,0.9,0.9}{
            \makecell[{{p{\linewidth}}}]{
                \texttt{\tiny{[P1$\langle$GM]}}
                \texttt{Straf: (1,2) is niet leeg, maar bevat 'C'.} \\
\\ 
\texttt{Jullie hebben gezamenlijk 4 van 8 straffen verzameld.} \\
\texttt{Probeer het opnieuw!} \\
            }
        }
    }
    & \\ \\

    \theutterance \stepcounter{utterance}  
    & \multicolumn{4}{p{0.6\linewidth}}{
        \cellcolor[rgb]{0.9,0.9,0.9}{
            \makecell[{{p{\linewidth}}}]{
                \texttt{\tiny{[P1$\rangle$GM]}}
                \texttt{Ik moet beter kijken naar het raster. Laat me de lege cellen identificeren. De lege cellen zijn gemarkeerd met '◌'. Mijn objecten staan op: C=(2,2), L=(1,1), P=(7,7). Doel: C=(4,4), L=(2,2), P=(6,6).} \\
\\ 
\texttt{Ik moet C van (2,2) naar (4,4) verplaatsen. Het raster heeft echter lijnen en randen die ik moet vermijden. Laat me eerst C verplaatsen naar een aangrenzende lege cel. Welke cellen grenzen aan (2,2)?} \\
\\ 
\texttt{Boven: (2,1) {-} bevat '╤' {-} niet leeg} \\
\texttt{Onder: (2,3) {-} bevat '┤' {-} niet leeg} \\
\texttt{Links: (1,2) {-} bevat 'C' (eigen positie?) Nee, dat is de cel zelf. Eigenlijk is (1,2) links van (2,2) {-} die bevat een lijn of letter.} \\
\texttt{Rechts: (3,2) {-} wat staat daar?} \\
\\ 
\texttt{Ik moet het raster zorgvuldig lezen. De '◌' symbolen geven lege cellen aan. Laten we die vinden:} \\
\\ 
\texttt{Rij 1: (4,1), (5,1), (6,1), (7,1) {-} maar wacht, er staat L op (1,1) en (2,1) en (3,1) hebben lijnen? Het raster toont L│◌│◌◌◌ dus L op 1,1, dan lijn, dan ◌ op 3,1? Nee, laten we het raster opnieuw interpreteren.} \\
\\ 
\texttt{```} \\
\texttt{ 1234567} \\
\texttt{╔═╤═╤═══╗} \\
\texttt{║L│◌│◌◌◌║ 1} \\
\texttt{║C├─┤◌◌◌║ 2} \\
\texttt{╟─┤◌├─┬─╢ 3} \\
\texttt{║◌│◌│◌│◌║ 4} \\
\texttt{╟─┼─┘◌│◌║ 5} \\
\texttt{║◌│◌◌┌┴─╢ 6} \\
\texttt{║◌│◌◌│◌P║ 7} \\
\texttt{╚═╧══╧══╝} \\
\texttt{```} \\
\\ 
\texttt{Kolommen 1{-}7, rijen 1{-}7.} \\
\\ 
\texttt{Laten we per cel kijken wat er staat (alleen de inhoud binnen de rasterlijnen):} \\
\\ 
\texttt{Rij 1: (1,1)=L, (2,1)=│, (3,1)=◌, (4,1)=◌, (5,1)=◌, (6,1)=◌, (7,1)=◌? Nee, het patroon is "║L│◌│◌◌◌║" dus na de linkerrand: L, │, ◌, │, ◌, ◌, ◌, rechterrand. Dus:} \\
\texttt{(1,1)=L, (2,1)=│, (3,1)=◌, (4,1)=│? Wacht, het format is lastig. "║" = rand, dan "L", dan "│", dan "◌", dan "│"? Nee, "║L│◌│◌◌◌║" {-} dit zijn 7 karakters tussen de randen: L, │, ◌, │, ◌, ◌, ◌. Dus:} \\
\texttt{(1,1)=L} \\
\texttt{(2,1)=│ (verticale lijn)} \\
\texttt{(3,1)=◌} \\
\texttt{(4,1)=│ (verticale lijn)} \\
\texttt{(5,1)=◌} \\
\texttt{(6,1)=◌} \\
\texttt{(7,1)=◌} \\
\\ 
\texttt{Rij 2: "║C├─┤◌◌◌║" {-}> C, ├, ─, ┤, ◌, ◌, ◌} \\
\texttt{(1,2)=C} \\
\texttt{(2,2)=├} \\
\texttt{(3,2)=─} \\
\texttt{(4,2)=┤} \\
\texttt{(5,2)=◌} \\
\texttt{(6,2)=◌} \\
\texttt{(7,2)=◌} \\
\\ 
\texttt{Rij 3: "╟─┤◌├─┬─╢" {-}> dit is een randrij? "╟" en "╢" zijn randstukken. Inhoud: ─, ┤, ◌, ├, ─, ┬, ─? Nee, "╟─┤◌├─┬─╢" tussen de randen: ─, ┤, ◌, ├, ─, ┬, ─. Dit is grotendeels een lijnrij.} \\
\\ 
\texttt{Rij 4: "║◌│◌│◌│◌║" {-}> ◌, │, ◌, │, ◌, │, ◌} \\
\texttt{(1,4)=◌} \\
\texttt{(2,4)=│} \\
\texttt{(3,4)=◌} \\
\texttt{(4,4)=│} \\
\texttt{(5,4)=◌} \\
\texttt{(6,4)=│} \\
\texttt{(7,4)=◌} \\
\\ 
\texttt{Rij 5: "╟─┼─┘◌│◌║" {-}> ─, ┼, ─, ┘, ◌, │, ◌} \\
\texttt{(1,5)=─} \\
\texttt{(2,5)=┼} \\
\texttt{(3,5)=─} \\
\texttt{(4,5)=┘} \\
\texttt{(5,5)=◌} \\
\texttt{(6,5)=│} \\
\texttt{(7,5)=◌} \\
\\ 
\texttt{Rij 6: "║◌│◌◌┌┴─╢" {-}> ◌, │, ◌, ◌, ┌, ┴, ─} \\
\texttt{(1,6)=◌} \\
\texttt{(2,6)=│} \\
\texttt{(3,6)=◌} \\
\texttt{(4,6)=◌} \\
\texttt{(5,6)=┌} \\
\texttt{(6,6)=┴} \\
\texttt{(7,6)=─} \\
\\ 
\texttt{Rij 7: "║◌│◌◌│◌P║" {-}> ◌, │, ◌, ◌, │, ◌, P} \\
\texttt{(1,7)=◌} \\
\texttt{(2,7)=│} \\
\texttt{(3,7)=◌} \\
\texttt{(4,7)=◌} \\
\texttt{(5,7)=│} \\
\texttt{(6,7)=◌} \\
\texttt{(7,7)=P} \\
\\ 
\texttt{Maar wacht, de regels zeggen: objecten staan op cellen. De objecten L, C, P staan volgens de beschrijving op het raster. Maar de coördinaten die ik eerder gaf (C=(2,2), L=(1,1), P=(7,7)) {-} dat klopt niet met wat ik nu zie. C staat op (1,2) in mijn analyse? In rij 2: C is het eerste karakter, dus (1,2). L staat op (1,1). P staat op (7,7).} \\
\\ 
\texttt{Dus mijn objecten zijn:} \\
\texttt{C = (1,2)} \\
\texttt{L = (1,1)} \\
\texttt{P = (7,7)} \\
\\ 
\texttt{Maar de lege cellen (◌) zijn:} \\
\texttt{(3,1), (5,1), (6,1), (7,1)} \\
\texttt{(5,2), (6,2), (7,2)} \\
\texttt{(3,3)? Rij 3 is "╟─┤◌├─┬─╢" dus (3,3)=◌} \\
\texttt{(1,4), (3,4), (5,4), (7,4)} \\
\texttt{(5,5), (7,5)} \\
\texttt{(1,6), (3,6), (4,6)} \\
\texttt{(1,7), (3,7), (4,7), (6,7)} \\
\\ 
\texttt{Doel: C=(4,4) {-} maar (4,4) bevat │, niet leeg! L=(2,2) bevat ├, niet leeg! P=(6,6) bevat ┴, niet leeg!} \\
\\ 
\texttt{Dit is een probleem. Het raster zoals getekend toont dat (4,4), (2,2), en (6,6) ALLEMAAL lijnstukken bevatten, geen lege cellen. We moeten onze doelconfiguratie herzien.} \\
\\ 
\texttt{ZEG: We hebben een probleem! Ik heb het raster verkeerd geïnterpreteerd. (4,4), (2,2), en (6,6) bevatten lijnstukken, geen lege cellen. Mijn objecten staan op: C=(1,2), L=(1,1), P=(7,7). Lege cellen die ik kan zien: (3,1),(5,1),(6,1),(7,1),(5,2),(6,2),(7,2),(3,3),(1,4),(3,4),(5,4),(7,4),(5,5),(7,5),(1,6),(3,6),(4,6),(1,7),(3,7),(4,7),(6,7). Welke lege cellen heb jij? We moeten nieuwe doelcoördinaten kiezen die voor ons beiden leeg zijn.} \\
            }
        }
    }
    & & \\ \\

    \theutterance \stepcounter{utterance}  
    & & & \multicolumn{2}{p{0.3\linewidth}}{
        \cellcolor[rgb]{0.9,0.9,0.9}{
            \makecell[{{p{\linewidth}}}]{
                \texttt{\tiny{[GM$|$GM]}}
                \texttt{Ik moet beter kijken naar het raster. Laat me de lege cellen identificeren. De lege cellen zijn gemarkeerd met '◌'. Mijn objecten staan op: C=(2,2), L=(1,1), P=(7,7). Doel: C=(4,4), L=(2,2), P=(6,6).} \\
\\ 
\texttt{Ik moet C van (2,2) naar (4,4) verplaatsen. Het raster heeft echter lijnen en randen die ik moet vermijden. Laat me eerst C verplaatsen naar een aangrenzende lege cel. Welke cellen grenzen aan (2,2)?} \\
\\ 
\texttt{Boven: (2,1) {-} bevat '╤' {-} niet leeg} \\
\texttt{Onder: (2,3) {-} bevat '┤' {-} niet leeg} \\
\texttt{Links: (1,2) {-} bevat 'C' (eigen positie?) Nee, dat is de cel zelf. Eigenlijk is (1,2) links van (2,2) {-} die bevat een lijn of letter.} \\
\texttt{Rechts: (3,2) {-} wat staat daar?} \\
\\ 
\texttt{Ik moet het raster zorgvuldig lezen. De '◌' symbolen geven lege cellen aan. Laten we die vinden:} \\
\\ 
\texttt{Rij 1: (4,1), (5,1), (6,1), (7,1) {-} maar wacht, er staat L op (1,1) en (2,1) en (3,1) hebben lijnen? Het raster toont L│◌│◌◌◌ dus L op 1,1, dan lijn, dan ◌ op 3,1? Nee, laten we het raster opnieuw interpreteren.} \\
\\ 
\texttt{```} \\
\texttt{ 1234567} \\
\texttt{╔═╤═╤═══╗} \\
\texttt{║L│◌│◌◌◌║ 1} \\
\texttt{║C├─┤◌◌◌║ 2} \\
\texttt{╟─┤◌├─┬─╢ 3} \\
\texttt{║◌│◌│◌│◌║ 4} \\
\texttt{╟─┼─┘◌│◌║ 5} \\
\texttt{║◌│◌◌┌┴─╢ 6} \\
\texttt{║◌│◌◌│◌P║ 7} \\
\texttt{╚═╧══╧══╝} \\
\texttt{```} \\
\\ 
\texttt{Kolommen 1{-}7, rijen 1{-}7.} \\
\\ 
\texttt{Laten we per cel kijken wat er staat (alleen de inhoud binnen de rasterlijnen):} \\
\\ 
\texttt{Rij 1: (1,1)=L, (2,1)=│, (3,1)=◌, (4,1)=◌, (5,1)=◌, (6,1)=◌, (7,1)=◌? Nee, het patroon is "║L│◌│◌◌◌║" dus na de linkerrand: L, │, ◌, │, ◌, ◌, ◌, rechterrand. Dus:} \\
\texttt{(1,1)=L, (2,1)=│, (3,1)=◌, (4,1)=│? Wacht, het format is lastig. "║" = rand, dan "L", dan "│", dan "◌", dan "│"? Nee, "║L│◌│◌◌◌║" {-} dit zijn 7 karakters tussen de randen: L, │, ◌, │, ◌, ◌, ◌. Dus:} \\
\texttt{(1,1)=L} \\
\texttt{(2,1)=│ (verticale lijn)} \\
\texttt{(3,1)=◌} \\
\texttt{(4,1)=│ (verticale lijn)} \\
\texttt{(5,1)=◌} \\
\texttt{(6,1)=◌} \\
\texttt{(7,1)=◌} \\
\\ 
\texttt{Rij 2: "║C├─┤◌◌◌║" {-}> C, ├, ─, ┤, ◌, ◌, ◌} \\
\texttt{(1,2)=C} \\
\texttt{(2,2)=├} \\
\texttt{(3,2)=─} \\
\texttt{(4,2)=┤} \\
\texttt{(5,2)=◌} \\
\texttt{(6,2)=◌} \\
\texttt{(7,2)=◌} \\
\\ 
\texttt{Rij 3: "╟─┤◌├─┬─╢" {-}> dit is een randrij? "╟" en "╢" zijn randstukken. Inhoud: ─, ┤, ◌, ├, ─, ┬, ─? Nee, "╟─┤◌├─┬─╢" tussen de randen: ─, ┤, ◌, ├, ─, ┬, ─. Dit is grotendeels een lijnrij.} \\
\\ 
\texttt{Rij 4: "║◌│◌│◌│◌║" {-}> ◌, │, ◌, │, ◌, │, ◌} \\
\texttt{(1,4)=◌} \\
\texttt{(2,4)=│} \\
\texttt{(3,4)=◌} \\
\texttt{(4,4)=│} \\
\texttt{(5,4)=◌} \\
\texttt{(6,4)=│} \\
\texttt{(7,4)=◌} \\
\\ 
\texttt{Rij 5: "╟─┼─┘◌│◌║" {-}> ─, ┼, ─, ┘, ◌, │, ◌} \\
\texttt{(1,5)=─} \\
\texttt{(2,5)=┼} \\
\texttt{(3,5)=─} \\
\texttt{(4,5)=┘} \\
\texttt{(5,5)=◌} \\
\texttt{(6,5)=│} \\
\texttt{(7,5)=◌} \\
\\ 
\texttt{Rij 6: "║◌│◌◌┌┴─╢" {-}> ◌, │, ◌, ◌, ┌, ┴, ─} \\
\texttt{(1,6)=◌} \\
\texttt{(2,6)=│} \\
\texttt{(3,6)=◌} \\
\texttt{(4,6)=◌} \\
\texttt{(5,6)=┌} \\
\texttt{(6,6)=┴} \\
\texttt{(7,6)=─} \\
\\ 
\texttt{Rij 7: "║◌│◌◌│◌P║" {-}> ◌, │, ◌, ◌, │, ◌, P} \\
\texttt{(1,7)=◌} \\
\texttt{(2,7)=│} \\
\texttt{(3,7)=◌} \\
\texttt{(4,7)=◌} \\
\texttt{(5,7)=│} \\
\texttt{(6,7)=◌} \\
\texttt{(7,7)=P} \\
\\ 
\texttt{Maar wacht, de regels zeggen: objecten staan op cellen. De objecten L, C, P staan volgens de beschrijving op het raster. Maar de coördinaten die ik eerder gaf (C=(2,2), L=(1,1), P=(7,7)) {-} dat klopt niet met wat ik nu zie. C staat op (1,2) in mijn analyse? In rij 2: C is het eerste karakter, dus (1,2). L staat op (1,1). P staat op (7,7).} \\
\\ 
\texttt{Dus mijn objecten zijn:} \\
\texttt{C = (1,2)} \\
\texttt{L = (1,1)} \\
\texttt{P = (7,7)} \\
\\ 
\texttt{Maar de lege cellen (◌) zijn:} \\
\texttt{(3,1), (5,1), (6,1), (7,1)} \\
\texttt{(5,2), (6,2), (7,2)} \\
\texttt{(3,3)? Rij 3 is "╟─┤◌├─┬─╢" dus (3,3)=◌} \\
\texttt{(1,4), (3,4), (5,4), (7,4)} \\
\texttt{(5,5), (7,5)} \\
\texttt{(1,6), (3,6), (4,6)} \\
\texttt{(1,7), (3,7), (4,7), (6,7)} \\
\\ 
\texttt{Doel: C=(4,4) {-} maar (4,4) bevat │, niet leeg! L=(2,2) bevat ├, niet leeg! P=(6,6) bevat ┴, niet leeg!} \\
\\ 
\texttt{Dit is een probleem. Het raster zoals getekend toont dat (4,4), (2,2), en (6,6) ALLEMAAL lijnstukken bevatten, geen lege cellen. We moeten onze doelconfiguratie herzien.} \\
\\ 
\texttt{ZEG: We hebben een probleem! Ik heb het raster verkeerd geïnterpreteerd. (4,4), (2,2), en (6,6) bevatten lijnstukken, geen lege cellen. Mijn objecten staan op: C=(1,2), L=(1,1), P=(7,7). Lege cellen die ik kan zien: (3,1),(5,1),(6,1),(7,1),(5,2),(6,2),(7,2),(3,3),(1,4),(3,4),(5,4),(7,4),(5,5),(7,5),(1,6),(3,6),(4,6),(1,7),(3,7),(4,7),(6,7). Welke lege cellen heb jij? We moeten nieuwe doelcoördinaten kiezen die voor ons beiden leeg zijn.} \\
            }
        }
    }
    & & \\ \\

    \theutterance \stepcounter{utterance}  
    & & & \multicolumn{2}{p{0.3\linewidth}}{
        \cellcolor[rgb]{0.9,0.9,0.9}{
            \makecell[{{p{\linewidth}}}]{
                \texttt{\tiny{[GM$|$GM]}}
                \texttt{Invalid format: head is not empty: 'Ik moet beter kijken naar het raster. Laat me de lege cellen identificeren. De lege cellen zijn gemarkeerd met '◌'. Mijn objecten staan op: C=(2,2), L=(1,1), P=(7,7). Doel: C=(4,4), L=(2,2), P=(6,6).} \\
\\ 
\texttt{Ik moet C van (2,2) naar (4,4) verplaatsen. Het raster heeft echter lijnen en randen die ik moet vermijden. Laat me eerst C verplaatsen naar een aangrenzende lege cel. Welke cellen grenzen aan (2,2)?} \\
\\ 
\texttt{Boven: (2,1) {-} bevat '╤' {-} niet leeg} \\
\texttt{Onder: (2,3) {-} bevat '┤' {-} niet leeg} \\
\texttt{Links: (1,2) {-} bevat 'C' (eigen positie?) Nee, dat is de cel zelf. Eigenlijk is (1,2) links van (2,2) {-} die bevat een lijn of letter.} \\
\texttt{Rechts: (3,2) {-} wat staat daar?} \\
\\ 
\texttt{Ik moet het raster zorgvuldig lezen. De '◌' symbolen geven lege cellen aan. Laten we die vinden:} \\
\\ 
\texttt{Rij 1: (4,1), (5,1), (6,1), (7,1) {-} maar wacht, er staat L op (1,1) en (2,1) en (3,1) hebben lijnen? Het raster toont L│◌│◌◌◌ dus L op 1,1, dan lijn, dan ◌ op 3,1? Nee, laten we het raster opnieuw interpreteren.} \\
\\ 
\\ 
\texttt{ 1234567} \\
\texttt{╔═╤═╤═══╗} \\
\texttt{║L│◌│◌◌◌║ 1} \\
\texttt{║C├─┤◌◌◌║ 2} \\
\texttt{╟─┤◌├─┬─╢ 3} \\
\texttt{║◌│◌│◌│◌║ 4} \\
\texttt{╟─┼─┘◌│◌║ 5} \\
\texttt{║◌│◌◌┌┴─╢ 6} \\
\texttt{║◌│◌◌│◌P║ 7} \\
\texttt{╚═╧══╧══╝} \\
\\ 
\\ 
\texttt{Kolommen 1{-}7, rijen 1{-}7.} \\
\\ 
\texttt{Laten we per cel kijken wat er staat (alleen de inhoud binnen de rasterlijnen):} \\
\\ 
\texttt{Rij 1: (1,1)=L, (2,1)=│, (3,1)=◌, (4,1)=◌, (5,1)=◌, (6,1)=◌, (7,1)=◌? Nee, het patroon is "║L│◌│◌◌◌║" dus na de linkerrand: L, │, ◌, │, ◌, ◌, ◌, rechterrand. Dus:} \\
\texttt{(1,1)=L, (2,1)=│, (3,1)=◌, (4,1)=│? Wacht, het format is lastig. "║" = rand, dan "L", dan "│", dan "◌", dan "│"? Nee, "║L│◌│◌◌◌║" {-} dit zijn 7 karakters tussen de randen: L, │, ◌, │, ◌, ◌, ◌. Dus:} \\
\texttt{(1,1)=L} \\
\texttt{(2,1)=│ (verticale lijn)} \\
\texttt{(3,1)=◌} \\
\texttt{(4,1)=│ (verticale lijn)} \\
\texttt{(5,1)=◌} \\
\texttt{(6,1)=◌} \\
\texttt{(7,1)=◌} \\
\\ 
\texttt{Rij 2: "║C├─┤◌◌◌║" {-}> C, ├, ─, ┤, ◌, ◌, ◌} \\
\texttt{(1,2)=C} \\
\texttt{(2,2)=├} \\
\texttt{(3,2)=─} \\
\texttt{(4,2)=┤} \\
\texttt{(5,2)=◌} \\
\texttt{(6,2)=◌} \\
\texttt{(7,2)=◌} \\
\\ 
\texttt{Rij 3: "╟─┤◌├─┬─╢" {-}> dit is een randrij? "╟" en "╢" zijn randstukken. Inhoud: ─, ┤, ◌, ├, ─, ┬, ─? Nee, "╟─┤◌├─┬─╢" tussen de randen: ─, ┤, ◌, ├, ─, ┬, ─. Dit is grotendeels een lijnrij.} \\
\\ 
\texttt{Rij 4: "║◌│◌│◌│◌║" {-}> ◌, │, ◌, │, ◌, │, ◌} \\
\texttt{(1,4)=◌} \\
\texttt{(2,4)=│} \\
\texttt{(3,4)=◌} \\
\texttt{(4,4)=│} \\
\texttt{(5,4)=◌} \\
\texttt{(6,4)=│} \\
\texttt{(7,4)=◌} \\
\\ 
\texttt{Rij 5: "╟─┼─┘◌│◌║" {-}> ─, ┼, ─, ┘, ◌, │, ◌} \\
\texttt{(1,5)=─} \\
\texttt{(2,5)=┼} \\
\texttt{(3,5)=─} \\
\texttt{(4,5)=┘} \\
\texttt{(5,5)=◌} \\
\texttt{(6,5)=│} \\
\texttt{(7,5)=◌} \\
\\ 
\texttt{Rij 6: "║◌│◌◌┌┴─╢" {-}> ◌, │, ◌, ◌, ┌, ┴, ─} \\
\texttt{(1,6)=◌} \\
\texttt{(2,6)=│} \\
\texttt{(3,6)=◌} \\
\texttt{(4,6)=◌} \\
\texttt{(5,6)=┌} \\
\texttt{(6,6)=┴} \\
\texttt{(7,6)=─} \\
\\ 
\texttt{Rij 7: "║◌│◌◌│◌P║" {-}> ◌, │, ◌, ◌, │, ◌, P} \\
\texttt{(1,7)=◌} \\
\texttt{(2,7)=│} \\
\texttt{(3,7)=◌} \\
\texttt{(4,7)=◌} \\
\texttt{(5,7)=│} \\
\texttt{(6,7)=◌} \\
\texttt{(7,7)=P} \\
\\ 
\texttt{Maar wacht, de regels zeggen: objecten staan op cellen. De objecten L, C, P staan volgens de beschrijving op het raster. Maar de coördinaten die ik eerder gaf (C=(2,2), L=(1,1), P=(7,7)) {-} dat klopt niet met wat ik nu zie. C staat op (1,2) in mijn analyse? In rij 2: C is het eerste karakter, dus (1,2). L staat op (1,1). P staat op (7,7).} \\
\\ 
\texttt{Dus mijn objecten zijn:} \\
\texttt{C = (1,2)} \\
\texttt{L = (1,1)} \\
\texttt{P = (7,7)} \\
\\ 
\texttt{Maar de lege cellen (◌) zijn:} \\
\texttt{(3,1), (5,1), (6,1), (7,1)} \\
\texttt{(5,2), (6,2), (7,2)} \\
\texttt{(3,3)? Rij 3 is "╟─┤◌├─┬─╢" dus (3,3)=◌} \\
\texttt{(1,4), (3,4), (5,4), (7,4)} \\
\texttt{(5,5), (7,5)} \\
\texttt{(1,6), (3,6), (4,6)} \\
\texttt{(1,7), (3,7), (4,7), (6,7)} \\
\\ 
\texttt{Doel: C=(4,4) {-} maar (4,4) bevat │, niet leeg! L=(2,2) bevat ├, niet leeg! P=(6,6) bevat ┴, niet leeg!} \\
\\ 
\texttt{Dit is een probleem. Het raster zoals getekend toont dat (4,4), (2,2), en (6,6) ALLEMAAL lijnstukken bevatten, geen lege cellen. We moeten onze doelconfiguratie herzien.} \\
\\ 
\texttt{'} \\
            }
        }
    }
    & & \\ \\

    \theutterance \stepcounter{utterance}  
    & & \multicolumn{4}{p{0.6\linewidth}}{
        \cellcolor[rgb]{0.9,0.9,0.9}{
            \makecell[{{p{\linewidth}}}]{
                \texttt{\tiny{[P1$\langle$GM]}}
                \texttt{Straf: Je bericht mag niets vóór het commando bevatten!} \\
\\ 
\texttt{Zorg ervoor dat je antwoord alleen `ZEG: <BERICHT>` of `BEWEEG: <OBJECT>, (<X>, <Y>)` bevat, en niets anders!} \\
\\ 
\texttt{Jullie hebben gezamenlijk 5 van 8 straffen verzameld.} \\
\texttt{Probeer het opnieuw!} \\
            }
        }
    }
    & \\ \\

    \theutterance \stepcounter{utterance}  
    & \multicolumn{4}{p{0.6\linewidth}}{
        \cellcolor[rgb]{0.9,0.9,0.9}{
            \makecell[{{p{\linewidth}}}]{
                \texttt{\tiny{[P1$\rangle$GM]}}
                \texttt{ZEG: We hebben een probleem! Ik heb het raster verkeerd geinterpreteerd. (4,4), (2,2), en (6,6) bevatten lijnstukken, geen lege cellen. Mijn objecten staan op: C=(1,2), L=(1,1), P=(7,7). Lege cellen die ik zie: (3,1),(5,1),(6,1),(7,1),(5,2),(6,2),(7,2),(3,3),(1,4),(3,4),(5,4),(7,4),(5,5),(7,5),(1,6),(3,6),(4,6),(1,7),(3,7),(4,7),(6,7). Welke lege cellen heb jij? We moeten nieuwe doelcoordinaten kiezen die voor ons beiden leeg zijn.} \\
            }
        }
    }
    & & \\ \\

    \theutterance \stepcounter{utterance}  
    & & & \multicolumn{2}{p{0.3\linewidth}}{
        \cellcolor[rgb]{0.9,0.9,0.9}{
            \makecell[{{p{\linewidth}}}]{
                \texttt{\tiny{[GM$|$GM]}}
                \texttt{ZEG: We hebben een probleem! Ik heb het raster verkeerd geinterpreteerd. (4,4), (2,2), en (6,6) bevatten lijnstukken, geen lege cellen. Mijn objecten staan op: C=(1,2), L=(1,1), P=(7,7). Lege cellen die ik zie: (3,1),(5,1),(6,1),(7,1),(5,2),(6,2),(7,2),(3,3),(1,4),(3,4),(5,4),(7,4),(5,5),(7,5),(1,6),(3,6),(4,6),(1,7),(3,7),(4,7),(6,7). Welke lege cellen heb jij? We moeten nieuwe doelcoordinaten kiezen die voor ons beiden leeg zijn.} \\
            }
        }
    }
    & & \\ \\

    \theutterance \stepcounter{utterance}  
    & & \multicolumn{4}{p{0.6\linewidth}}{
        \cellcolor[rgb]{0.9,0.9,0.9}{
            \makecell[{{p{\linewidth}}}]{
                \texttt{\tiny{[GM$\rangle$P2]}}
                \texttt{Je bericht is doorgestuurd naar de andere speler.} \\
\\ 
\\ 
\texttt{Je speelt momenteel ronde 4 van maximaal 12.} \\
\texttt{Jullie hebben gezamenlijk 5 van 8 straffen verzameld.} \\
\texttt{De andere speler stuurde dit bericht:} \\
\\ 
\\ 
\texttt{"We hebben een probleem! Ik heb het raster verkeerd geinterpreteerd. (4,4), (2,2), en (6,6) bevatten lijnstukken, geen lege cellen. Mijn objecten staan op: C=(1,2), L=(1,1), P=(7,7). Lege cellen die ik zie: (3,1),(5,1),(6,1),(7,1),(5,2),(6,2),(7,2),(3,3),(1,4),(3,4),(5,4),(7,4),(5,5),(7,5),(1,6),(3,6),(4,6),(1,7),(3,7),(4,7),(6,7). Welke lege cellen heb jij? We moeten nieuwe doelcoordinaten kiezen die voor ons beiden leeg zijn."} \\
\\ 
\\ 
\texttt{Wat is je volgende opdracht?} \\
            }
        }
    }
    & \\ \\

    \theutterance \stepcounter{utterance}  
    & & & \multicolumn{4}{p{0.6\linewidth}}{
        \cellcolor[rgb]{0.9,0.9,0.9}{
            \makecell[{{p{\linewidth}}}]{
                \texttt{\tiny{[GM$\langle$P2]}}
                \texttt{Eerst ga ik mijn raster zorgvuldig bekijken om te bepalen welke cellen leeg zijn.} \\
\\ 
\texttt{Mijn raster:} \\
\texttt{```} \\
\texttt{ 1234567} \\
\texttt{╔═╤═╤═══╗} \\
\texttt{║◌│◌│◌◌◌║ 1} \\
\texttt{║◌├─┤◌L◌║ 2} \\
\texttt{╟─┤P├─┬─╢ 3} \\
\texttt{║◌│◌│◌│◌║ 4} \\
\texttt{╟─┼─┘C│◌║ 5} \\
\texttt{║◌│◌◌┌┴─╢ 6} \\
\texttt{║◌│◌◌│◌◌║ 7} \\
\texttt{╚═╧══╧══╝} \\
\texttt{```} \\
\\ 
\texttt{Lege cellen (◌) op mijn bord:} \\
\texttt{{-} Rij 1: (1,1), (2,1), (4,1), (5,1), (6,1), (7,1)} \\
\texttt{{-} Rij 2: (1,2), (4,2), (6,2), (7,2)} \\
\texttt{{-} Rij 3: (1,3), (4,3), (6,3), (7,3)} \\
\texttt{{-} Rij 4: (1,4), (3,4), (5,4), (7,4)} \\
\texttt{{-} Rij 5: (1,5), (2,5), (5,5), (7,5)} \\
\texttt{{-} Rij 6: (1,6), (2,6), (5,6), (6,6)} \\
\texttt{{-} Rij 7: (1,7), (2,7), (4,7), (5,7), (6,7), (7,7)} \\
\\ 
\texttt{Gemeenschappelijke lege cellen (uit overlap met die van de ander):} \\
\texttt{(5,1), (6,1), (7,1), (6,2), (7,2), (7,4), (5,5), (7,5), (1,6), (1,7), (4,7), (6,7)} \\
\\ 
\texttt{`ZEG: Mijn lege cellen: (1,1),(2,1),(4,1),(5,1),(6,1),(7,1),(1,2),(4,2),(6,2),(7,2),(1,3),(4,3),(6,3),(7,3),(1,4),(3,4),(5,4),(7,4),(1,5),(2,5),(5,5),(7,5),(1,6),(2,6),(5,6),(6,6),(1,7),(2,7),(4,7),(5,7),(6,7),(7,7). Overlap tussen onze lege cellen: (5,1),(6,1),(7,1),(6,2),(7,2),(7,4),(5,5),(7,5),(1,6),(1,7),(4,7),(6,7). Laten we hieruit drie kiezen, bijvoorbeeld C=(5,5), L=(6,2), P=(7,4)?`} \\
            }
        }
    }
     \\ \\

    \theutterance \stepcounter{utterance}  
    & & & \multicolumn{2}{p{0.3\linewidth}}{
        \cellcolor[rgb]{0.9,0.9,0.9}{
            \makecell[{{p{\linewidth}}}]{
                \texttt{\tiny{[GM$|$GM]}}
                \texttt{Eerst ga ik mijn raster zorgvuldig bekijken om te bepalen welke cellen leeg zijn.} \\
\\ 
\texttt{Mijn raster:} \\
\texttt{```} \\
\texttt{ 1234567} \\
\texttt{╔═╤═╤═══╗} \\
\texttt{║◌│◌│◌◌◌║ 1} \\
\texttt{║◌├─┤◌L◌║ 2} \\
\texttt{╟─┤P├─┬─╢ 3} \\
\texttt{║◌│◌│◌│◌║ 4} \\
\texttt{╟─┼─┘C│◌║ 5} \\
\texttt{║◌│◌◌┌┴─╢ 6} \\
\texttt{║◌│◌◌│◌◌║ 7} \\
\texttt{╚═╧══╧══╝} \\
\texttt{```} \\
\\ 
\texttt{Lege cellen (◌) op mijn bord:} \\
\texttt{{-} Rij 1: (1,1), (2,1), (4,1), (5,1), (6,1), (7,1)} \\
\texttt{{-} Rij 2: (1,2), (4,2), (6,2), (7,2)} \\
\texttt{{-} Rij 3: (1,3), (4,3), (6,3), (7,3)} \\
\texttt{{-} Rij 4: (1,4), (3,4), (5,4), (7,4)} \\
\texttt{{-} Rij 5: (1,5), (2,5), (5,5), (7,5)} \\
\texttt{{-} Rij 6: (1,6), (2,6), (5,6), (6,6)} \\
\texttt{{-} Rij 7: (1,7), (2,7), (4,7), (5,7), (6,7), (7,7)} \\
\\ 
\texttt{Gemeenschappelijke lege cellen (uit overlap met die van de ander):} \\
\texttt{(5,1), (6,1), (7,1), (6,2), (7,2), (7,4), (5,5), (7,5), (1,6), (1,7), (4,7), (6,7)} \\
\\ 
\texttt{`ZEG: Mijn lege cellen: (1,1),(2,1),(4,1),(5,1),(6,1),(7,1),(1,2),(4,2),(6,2),(7,2),(1,3),(4,3),(6,3),(7,3),(1,4),(3,4),(5,4),(7,4),(1,5),(2,5),(5,5),(7,5),(1,6),(2,6),(5,6),(6,6),(1,7),(2,7),(4,7),(5,7),(6,7),(7,7). Overlap tussen onze lege cellen: (5,1),(6,1),(7,1),(6,2),(7,2),(7,4),(5,5),(7,5),(1,6),(1,7),(4,7),(6,7). Laten we hieruit drie kiezen, bijvoorbeeld C=(5,5), L=(6,2), P=(7,4)?`} \\
            }
        }
    }
    & & \\ \\

    \theutterance \stepcounter{utterance}  
    & & & \multicolumn{2}{p{0.3\linewidth}}{
        \cellcolor[rgb]{0.9,0.9,0.9}{
            \makecell[{{p{\linewidth}}}]{
                \texttt{\tiny{[GM$|$GM]}}
                \texttt{Invalid format: head is not empty: 'Eerst ga ik mijn raster zorgvuldig bekijken om te bepalen welke cellen leeg zijn.} \\
\\ 
\texttt{Mijn raster:} \\
\\ 
\texttt{ 1234567} \\
\texttt{╔═╤═╤═══╗} \\
\texttt{║◌│◌│◌◌◌║ 1} \\
\texttt{║◌├─┤◌L◌║ 2} \\
\texttt{╟─┤P├─┬─╢ 3} \\
\texttt{║◌│◌│◌│◌║ 4} \\
\texttt{╟─┼─┘C│◌║ 5} \\
\texttt{║◌│◌◌┌┴─╢ 6} \\
\texttt{║◌│◌◌│◌◌║ 7} \\
\texttt{╚═╧══╧══╝} \\
\\ 
\\ 
\texttt{Lege cellen (◌) op mijn bord:} \\
\texttt{{-} Rij 1: (1,1), (2,1), (4,1), (5,1), (6,1), (7,1)} \\
\texttt{{-} Rij 2: (1,2), (4,2), (6,2), (7,2)} \\
\texttt{{-} Rij 3: (1,3), (4,3), (6,3), (7,3)} \\
\texttt{{-} Rij 4: (1,4), (3,4), (5,4), (7,4)} \\
\texttt{{-} Rij 5: (1,5), (2,5), (5,5), (7,5)} \\
\texttt{{-} Rij 6: (1,6), (2,6), (5,6), (6,6)} \\
\texttt{{-} Rij 7: (1,7), (2,7), (4,7), (5,7), (6,7), (7,7)} \\
\\ 
\texttt{Gemeenschappelijke lege cellen (uit overlap met die van de ander):} \\
\texttt{(5,1), (6,1), (7,1), (6,2), (7,2), (7,4), (5,5), (7,5), (1,6), (1,7), (4,7), (6,7)} \\
\\ 
\texttt{'} \\
            }
        }
    }
    & & \\ \\

    \theutterance \stepcounter{utterance}  
    & & \multicolumn{4}{p{0.6\linewidth}}{
        \cellcolor[rgb]{0.9,0.9,0.9}{
            \makecell[{{p{\linewidth}}}]{
                \texttt{\tiny{[GM$\rangle$P2]}}
                \texttt{Straf: Je bericht mag niets vóór het commando bevatten!} \\
\\ 
\texttt{Zorg ervoor dat je antwoord alleen `ZEG: <BERICHT>` of `BEWEEG: <OBJECT>, (<X>, <Y>)` bevat, en niets anders!} \\
\\ 
\texttt{Jullie hebben gezamenlijk 6 van 8 straffen verzameld.} \\
\texttt{Probeer het opnieuw!} \\
            }
        }
    }
    & \\ \\

    \theutterance \stepcounter{utterance}  
    & & & \multicolumn{4}{p{0.6\linewidth}}{
        \cellcolor[rgb]{0.9,0.9,0.9}{
            \makecell[{{p{\linewidth}}}]{
                \texttt{\tiny{[GM$\langle$P2]}}
                \texttt{`ZEG: Mijn lege cellen: (1,1),(2,1),(4,1),(5,1),(6,1),(7,1),(1,2),(4,2),(6,2),(7,2),(1,3),(4,3),(6,3),(7,3),(1,4),(3,4),(5,4),(7,4),(1,5),(2,5),(5,5),(7,5),(1,6),(2,6),(5,6),(6,6),(1,7),(2,7),(4,7),(5,7),(6,7),(7,7). Overlap: (5,1),(6,1),(7,1),(6,2),(7,2),(7,4),(5,5),(7,5),(1,6),(1,7),(4,7),(6,7). Voorstel: C=(5,5), L=(6,2), P=(7,4)?`} \\
            }
        }
    }
     \\ \\

    \theutterance \stepcounter{utterance}  
    & & & \multicolumn{2}{p{0.3\linewidth}}{
        \cellcolor[rgb]{0.9,0.9,0.9}{
            \makecell[{{p{\linewidth}}}]{
                \texttt{\tiny{[GM$|$GM]}}
                \texttt{`ZEG: Mijn lege cellen: (1,1),(2,1),(4,1),(5,1),(6,1),(7,1),(1,2),(4,2),(6,2),(7,2),(1,3),(4,3),(6,3),(7,3),(1,4),(3,4),(5,4),(7,4),(1,5),(2,5),(5,5),(7,5),(1,6),(2,6),(5,6),(6,6),(1,7),(2,7),(4,7),(5,7),(6,7),(7,7). Overlap: (5,1),(6,1),(7,1),(6,2),(7,2),(7,4),(5,5),(7,5),(1,6),(1,7),(4,7),(6,7). Voorstel: C=(5,5), L=(6,2), P=(7,4)?`} \\
            }
        }
    }
    & & \\ \\

    \theutterance \stepcounter{utterance}  
    & & \multicolumn{4}{p{0.6\linewidth}}{
        \cellcolor[rgb]{0.9,0.9,0.9}{
            \makecell[{{p{\linewidth}}}]{
                \texttt{\tiny{[P1$\langle$GM]}}
                \texttt{Je bericht is doorgestuurd naar de andere speler.} \\
\\ 
\\ 
\texttt{Je speelt momenteel ronde 4 van maximaal 12.} \\
\texttt{Jullie hebben gezamenlijk 6 van 8 straffen verzameld.} \\
\texttt{De andere speler stuurde dit bericht:} \\
\\ 
\\ 
\texttt{"Mijn lege cellen: (1,1),(2,1),(4,1),(5,1),(6,1),(7,1),(1,2),(4,2),(6,2),(7,2),(1,3),(4,3),(6,3),(7,3),(1,4),(3,4),(5,4),(7,4),(1,5),(2,5),(5,5),(7,5),(1,6),(2,6),(5,6),(6,6),(1,7),(2,7),(4,7),(5,7),(6,7),(7,7). Overlap: (5,1),(6,1),(7,1),(6,2),(7,2),(7,4),(5,5),(7,5),(1,6),(1,7),(4,7),(6,7). Voorstel: C=(5,5), L=(6,2), P=(7,4)?"} \\
\\ 
\\ 
\texttt{Wat is je volgende opdracht?} \\
            }
        }
    }
    & \\ \\

    \theutterance \stepcounter{utterance}  
    & \multicolumn{4}{p{0.6\linewidth}}{
        \cellcolor[rgb]{0.9,0.9,0.9}{
            \makecell[{{p{\linewidth}}}]{
                \texttt{\tiny{[P1$\rangle$GM]}}
                \texttt{ZEG: Controleer of (5,5), (6,2) en (7,4) in mijn lijst staan. (5,5)=leeg ja, (6,2)=leeg ja, (7,4)=leeg ja. Akkoord met C=(5,5), L=(6,2), P=(7,4)! Ik begin met verplaatsen.} \\
            }
        }
    }
    & & \\ \\

    \theutterance \stepcounter{utterance}  
    & & & \multicolumn{2}{p{0.3\linewidth}}{
        \cellcolor[rgb]{0.9,0.9,0.9}{
            \makecell[{{p{\linewidth}}}]{
                \texttt{\tiny{[GM$|$GM]}}
                \texttt{ZEG: Controleer of (5,5), (6,2) en (7,4) in mijn lijst staan. (5,5)=leeg ja, (6,2)=leeg ja, (7,4)=leeg ja. Akkoord met C=(5,5), L=(6,2), P=(7,4)! Ik begin met verplaatsen.} \\
            }
        }
    }
    & & \\ \\

    \theutterance \stepcounter{utterance}  
    & & \multicolumn{4}{p{0.6\linewidth}}{
        \cellcolor[rgb]{0.9,0.9,0.9}{
            \makecell[{{p{\linewidth}}}]{
                \texttt{\tiny{[GM$\rangle$P2]}}
                \texttt{Je bericht is doorgestuurd naar de andere speler.} \\
\\ 
\\ 
\texttt{Je speelt momenteel ronde 5 van maximaal 12.} \\
\texttt{Jullie hebben gezamenlijk 6 van 8 straffen verzameld.} \\
\texttt{De andere speler stuurde dit bericht:} \\
\\ 
\\ 
\texttt{"Controleer of (5,5), (6,2) en (7,4) in mijn lijst staan. (5,5)=leeg ja, (6,2)=leeg ja, (7,4)=leeg ja. Akkoord met C=(5,5), L=(6,2), P=(7,4)! Ik begin met verplaatsen."} \\
\\ 
\\ 
\texttt{Wat is je volgende opdracht?} \\
            }
        }
    }
    & \\ \\

    \theutterance \stepcounter{utterance}  
    & & & \multicolumn{4}{p{0.6\linewidth}}{
        \cellcolor[rgb]{0.9,0.9,0.9}{
            \makecell[{{p{\linewidth}}}]{
                \texttt{\tiny{[GM$\langle$P2]}}
                \texttt{`BEWEEG: C, (5,5)`} \\
            }
        }
    }
     \\ \\

    \theutterance \stepcounter{utterance}  
    & & & \multicolumn{2}{p{0.3\linewidth}}{
        \cellcolor[rgb]{0.9,0.9,0.9}{
            \makecell[{{p{\linewidth}}}]{
                \texttt{\tiny{[GM$|$GM]}}
                \texttt{`BEWEEG: C, (5,5)`} \\
            }
        }
    }
    & & \\ \\

    \theutterance \stepcounter{utterance}  
    & & & \multicolumn{2}{p{0.3\linewidth}}{
        \cellcolor[rgb]{0.9,0.9,0.9}{
            \makecell[{{p{\linewidth}}}]{
                \texttt{\tiny{[GM$|$GM]}}
                \texttt{Straf: (5,5) is niet leeg, maar bevat 'C'.} \\
\\ 
\texttt{Jullie hebben gezamenlijk 7 van 8 straffen verzameld.} \\
\texttt{Probeer het opnieuw!} \\
            }
        }
    }
    & & \\ \\

    \theutterance \stepcounter{utterance}  
    & & \multicolumn{4}{p{0.6\linewidth}}{
        \cellcolor[rgb]{0.9,0.9,0.9}{
            \makecell[{{p{\linewidth}}}]{
                \texttt{\tiny{[GM$\rangle$P2]}}
                \texttt{Straf: (5,5) is niet leeg, maar bevat 'C'.} \\
\\ 
\texttt{Jullie hebben gezamenlijk 7 van 8 straffen verzameld.} \\
\texttt{Probeer het opnieuw!} \\
            }
        }
    }
    & \\ \\

    \theutterance \stepcounter{utterance}  
    & & & \multicolumn{4}{p{0.6\linewidth}}{
        \cellcolor[rgb]{0.9,0.9,0.9}{
            \makecell[{{p{\linewidth}}}]{
                \texttt{\tiny{[GM$\langle$P2]}}
                \texttt{Mijn C staat al op (4,5), dus ik moet controleren of (5,5) leeg is. Op mijn raster staat ◌ op (5,5) volgens mijn lijst, maar blijkbaar is dat niet zo. Laat ik mijn raster opnieuw bekijken.} \\
\\ 
\texttt{Mijn raster nogmaals:} \\
\texttt{Rij 5: `╟─┼─┘C│◌║ 5` dus kolom 5 is ◌.} \\
\\ 
\texttt{Mijn objecten: C=(4,5), L=(5,2), P=(3,3). Doel: C=(5,5), L=(6,2), P=(7,4).} \\
\\ 
\texttt{Voor C: van (4,5) naar (5,5). (5,5) is rij 5, kolom 5 {-} dat is het ◌ symbool. Maar de straf zegt dat (5,5) 'C' bevat. Dat is vreemd, want C staat op (4,5). Wacht, de strafmelding zegt dat (5,5) niet leeg is maar 'C' bevat {-} misschien interpreteert het systeem mijn eigen C als bezet. Maar ik probeer C te verplaatsen VAN (4,5) NAAR (5,5). Dat zou moeten kunnen als (5,5) leeg is.} \\
\\ 
\texttt{Misschien is (5,5) wel bezet door iets anders? Even opnieuw kijken naar rij 5: `╟─┼─┘C│◌║ 5`. De karakters: kolom 1=╟, 2=─, 3=┼, 4=─, 5=┘, 6=C, 7=│, dan ◌? Nee, wacht {-} de coördinaten lopen van 1 t/m 7. Het raster:} \\
\\ 
\texttt{Kolom 1,2,3,4,5,6,7 op rij 5: ╟,─,┼,─,┘,C,│,◌? } \\
\\ 
\texttt{Ik moet het raster beter lezen. Het lijkt erop dat de kolommen verspringen door de grafische tekens. C staat volgens mij op (4,5) maar is dat wel kolom 4?} \\
\\ 
\texttt{`╟─┼─┘C│◌║ 5` {-} positie 1=╟, 2=─, 3=┼, 4=─, 5=┘, 6=C, 7=│... waar is ◌? Het bord heeft 7 kolommen. Misschien is C op (6,5) en niet (4,5)?} \\
\\ 
\texttt{Laat ik het hele raster opnieuw analyseren met coördinaten.} \\
\\ 
\texttt{`BEWEEG: C, (6,2)`} \\
            }
        }
    }
     \\ \\

    \theutterance \stepcounter{utterance}  
    & & & \multicolumn{2}{p{0.3\linewidth}}{
        \cellcolor[rgb]{0.9,0.9,0.9}{
            \makecell[{{p{\linewidth}}}]{
                \texttt{\tiny{[GM$|$GM]}}
                \texttt{Mijn C staat al op (4,5), dus ik moet controleren of (5,5) leeg is. Op mijn raster staat ◌ op (5,5) volgens mijn lijst, maar blijkbaar is dat niet zo. Laat ik mijn raster opnieuw bekijken.} \\
\\ 
\texttt{Mijn raster nogmaals:} \\
\texttt{Rij 5: `╟─┼─┘C│◌║ 5` dus kolom 5 is ◌.} \\
\\ 
\texttt{Mijn objecten: C=(4,5), L=(5,2), P=(3,3). Doel: C=(5,5), L=(6,2), P=(7,4).} \\
\\ 
\texttt{Voor C: van (4,5) naar (5,5). (5,5) is rij 5, kolom 5 {-} dat is het ◌ symbool. Maar de straf zegt dat (5,5) 'C' bevat. Dat is vreemd, want C staat op (4,5). Wacht, de strafmelding zegt dat (5,5) niet leeg is maar 'C' bevat {-} misschien interpreteert het systeem mijn eigen C als bezet. Maar ik probeer C te verplaatsen VAN (4,5) NAAR (5,5). Dat zou moeten kunnen als (5,5) leeg is.} \\
\\ 
\texttt{Misschien is (5,5) wel bezet door iets anders? Even opnieuw kijken naar rij 5: `╟─┼─┘C│◌║ 5`. De karakters: kolom 1=╟, 2=─, 3=┼, 4=─, 5=┘, 6=C, 7=│, dan ◌? Nee, wacht {-} de coördinaten lopen van 1 t/m 7. Het raster:} \\
\\ 
\texttt{Kolom 1,2,3,4,5,6,7 op rij 5: ╟,─,┼,─,┘,C,│,◌? } \\
\\ 
\texttt{Ik moet het raster beter lezen. Het lijkt erop dat de kolommen verspringen door de grafische tekens. C staat volgens mij op (4,5) maar is dat wel kolom 4?} \\
\\ 
\texttt{`╟─┼─┘C│◌║ 5` {-} positie 1=╟, 2=─, 3=┼, 4=─, 5=┘, 6=C, 7=│... waar is ◌? Het bord heeft 7 kolommen. Misschien is C op (6,5) en niet (4,5)?} \\
\\ 
\texttt{Laat ik het hele raster opnieuw analyseren met coördinaten.} \\
\\ 
\texttt{`BEWEEG: C, (6,2)`} \\
            }
        }
    }
    & & \\ \\

    \theutterance \stepcounter{utterance}  
    & & & \multicolumn{2}{p{0.3\linewidth}}{
        \cellcolor[rgb]{0.9,0.9,0.9}{
            \makecell[{{p{\linewidth}}}]{
                \texttt{\tiny{[GM$|$GM]}}
                \texttt{Invalid format: head is not empty: 'Mijn C staat al op (4,5), dus ik moet controleren of (5,5) leeg is. Op mijn raster staat ◌ op (5,5) volgens mijn lijst, maar blijkbaar is dat niet zo. Laat ik mijn raster opnieuw bekijken.} \\
\\ 
\texttt{Mijn raster nogmaals:} \\
\texttt{Rij 5: ╟─┼─┘C│◌║ 5 dus kolom 5 is ◌.} \\
\\ 
\texttt{Mijn objecten: C=(4,5), L=(5,2), P=(3,3). Doel: C=(5,5), L=(6,2), P=(7,4).} \\
\\ 
\texttt{Voor C: van (4,5) naar (5,5). (5,5) is rij 5, kolom 5 {-} dat is het ◌ symbool. Maar de straf zegt dat (5,5) 'C' bevat. Dat is vreemd, want C staat op (4,5). Wacht, de strafmelding zegt dat (5,5) niet leeg is maar 'C' bevat {-} misschien interpreteert het systeem mijn eigen C als bezet. Maar ik probeer C te verplaatsen VAN (4,5) NAAR (5,5). Dat zou moeten kunnen als (5,5) leeg is.} \\
\\ 
\texttt{Misschien is (5,5) wel bezet door iets anders? Even opnieuw kijken naar rij 5: ╟─┼─┘C│◌║ 5. De karakters: kolom 1=╟, 2=─, 3=┼, 4=─, 5=┘, 6=C, 7=│, dan ◌? Nee, wacht {-} de coördinaten lopen van 1 t/m 7. Het raster:} \\
\\ 
\texttt{Kolom 1,2,3,4,5,6,7 op rij 5: ╟,─,┼,─,┘,C,│,◌? } \\
\\ 
\texttt{Ik moet het raster beter lezen. Het lijkt erop dat de kolommen verspringen door de grafische tekens. C staat volgens mij op (4,5) maar is dat wel kolom 4?} \\
\\ 
\texttt{╟─┼─┘C│◌║ 5 {-} positie 1=╟, 2=─, 3=┼, 4=─, 5=┘, 6=C, 7=│... waar is ◌? Het bord heeft 7 kolommen. Misschien is C op (6,5) en niet (4,5)?} \\
\\ 
\texttt{Laat ik het hele raster opnieuw analyseren met coördinaten.} \\
\\ 
\texttt{'} \\
            }
        }
    }
    & & \\ \\

    \theutterance \stepcounter{utterance}  
    & & \multicolumn{4}{p{0.6\linewidth}}{
        \cellcolor[rgb]{0.9,0.9,0.9}{
            \makecell[{{p{\linewidth}}}]{
                \texttt{\tiny{[GM$\rangle$P2]}}
                \texttt{Straf: Je bericht mag niets vóór het commando bevatten!} \\
\\ 
\texttt{Zorg ervoor dat je antwoord alleen `ZEG: <BERICHT>` of `BEWEEG: <OBJECT>, (<X>, <Y>)` bevat, en niets anders!} \\
\\ 
\texttt{Jullie hebben gezamenlijk 8 van 8 straffen verzameld.} \\
\texttt{Probeer het opnieuw!} \\
            }
        }
    }
    & \\ \\

    \theutterance \stepcounter{utterance}  
    & & & \multicolumn{4}{p{0.6\linewidth}}{
        \cellcolor[rgb]{0.9,0.9,0.9}{
            \makecell[{{p{\linewidth}}}]{
                \texttt{\tiny{[GM$\langle$P2]}}
                \texttt{We hebben 8 strafpunten. Het spel is verloren. Mijn excuses voor de fouten.} \\
            }
        }
    }
     \\ \\

    \theutterance \stepcounter{utterance}  
    & & & \multicolumn{2}{p{0.3\linewidth}}{
        \cellcolor[rgb]{0.9,0.9,0.9}{
            \makecell[{{p{\linewidth}}}]{
                \texttt{\tiny{[GM$|$GM]}}
                \texttt{We hebben 8 strafpunten. Het spel is verloren. Mijn excuses voor de fouten.} \\
            }
        }
    }
    & & \\ \\

    \theutterance \stepcounter{utterance}  
    & & & \multicolumn{2}{p{0.3\linewidth}}{
        \cellcolor[rgb]{0.9,0.9,0.9}{
            \makecell[{{p{\linewidth}}}]{
                \texttt{\tiny{[GM$|$GM]}}
                \texttt{Invalid response format} \\
            }
        }
    }
    & & \\ \\

    \theutterance \stepcounter{utterance}  
    & & & \multicolumn{2}{p{0.3\linewidth}}{
        \cellcolor[rgb]{0.9,0.9,0.9}{
            \makecell[{{p{\linewidth}}}]{
                \texttt{\tiny{[GM$|$GM]}}
                \texttt{Maximum number of penalties exceeded} \\
            }
        }
    }
    & & \\ \\

    \theutterance \stepcounter{utterance}  
    & & & \multicolumn{2}{p{0.3\linewidth}}{
        \cellcolor[rgb]{0.9,0.9,0.9}{
            \makecell[{{p{\linewidth}}}]{
                \texttt{\tiny{[GM$|$GM]}}
                \texttt{Initial states:} \\
\texttt{```} \\
\texttt{Player 1:} \\
\texttt{ 1234567} \\
\texttt{╔═╤═╤═══╗} \\
\texttt{║L│◌│◌◌◌║ 1} \\
\texttt{║C├─┤◌◌◌║ 2} \\
\texttt{╟─┤◌├─┬─╢ 3} \\
\texttt{║◌│◌│◌│◌║ 4} \\
\texttt{╟─┼─┘◌│◌║ 5} \\
\texttt{║◌│◌◌┌┴─╢ 6} \\
\texttt{║◌│◌◌│◌P║ 7} \\
\texttt{╚═╧══╧══╝} \\
\\ 
\\ 
\texttt{Player 2:} \\
\texttt{ 1234567} \\
\texttt{╔═╤═╤═══╗} \\
\texttt{║◌│◌│◌◌◌║ 1} \\
\texttt{║◌├─┤◌L◌║ 2} \\
\texttt{╟─┤P├─┬─╢ 3} \\
\texttt{║◌│◌│◌│◌║ 4} \\
\texttt{╟─┼─┘C│◌║ 5} \\
\texttt{║◌│◌◌┌┴─╢ 6} \\
\texttt{║◌│◌◌│◌◌║ 7} \\
\texttt{╚═╧══╧══╝} \\
\\ 
\texttt{```} \\
            }
        }
    }
    & & \\ \\

    \theutterance \stepcounter{utterance}  
    & & & \multicolumn{2}{p{0.3\linewidth}}{
        \cellcolor[rgb]{0.9,0.9,0.9}{
            \makecell[{{p{\linewidth}}}]{
                \texttt{\tiny{[GM$|$GM]}}
                \texttt{End states:} \\
\texttt{```} \\
\texttt{Player 1:} \\
\texttt{ 1234567} \\
\texttt{╔═╤═╤═══╗} \\
\texttt{║L│◌│◌◌◌║ 1} \\
\texttt{║C├─┤◌◌◌║ 2} \\
\texttt{╟─┤◌├─┬─╢ 3} \\
\texttt{║◌│◌│◌│◌║ 4} \\
\texttt{╟─┼─┘◌│◌║ 5} \\
\texttt{║◌│◌◌┌┴─╢ 6} \\
\texttt{║◌│◌◌│◌P║ 7} \\
\texttt{╚═╧══╧══╝} \\
\\ 
\\ 
\texttt{Player 2:} \\
\texttt{ 1234567} \\
\texttt{╔═╤═╤═══╗} \\
\texttt{║◌│◌│◌◌◌║ 1} \\
\texttt{║◌├─┤◌L◌║ 2} \\
\texttt{╟─┤P├─┬─╢ 3} \\
\texttt{║◌│◌│◌│◌║ 4} \\
\texttt{╟─┼─┘C│◌║ 5} \\
\texttt{║◌│◌◌┌┴─╢ 6} \\
\texttt{║◌│◌◌│◌◌║ 7} \\
\texttt{╚═╧══╧══╝} \\
\\ 
\texttt{```} \\
            }
        }
    }
    & & \\ \\

    \theutterance \stepcounter{utterance}  
    & & & \multicolumn{2}{p{0.3\linewidth}}{
        \cellcolor[rgb]{0.9,0.9,0.9}{
            \makecell[{{p{\linewidth}}}]{
                \texttt{\tiny{[GM$|$GM]}}
                \texttt{* success: False} \\
\texttt{* lose: False} \\
\texttt{* aborted: True} \\
\texttt{{-}{-}{-}{-}{-}{-}{-}} \\
\texttt{* Shifts: 0.00} \\
\texttt{* Max Shifts: 4.00} \\
\texttt{* Min Shifts: 2.00} \\
\texttt{* End Distance Sum: 15.76} \\
\texttt{* Init Distance Sum: 15.76} \\
\texttt{* Expected Distance Sum: 12.57} \\
\texttt{* Penalties: 9.00} \\
\texttt{* Max Penalties: 8.00} \\
\texttt{* Rounds: 4.00} \\
\texttt{* Max Rounds: 12.00} \\
\texttt{* Object Count: 3.00} \\
            }
        }
    }
    & & \\ \\

\end{supertabular}
}

\end{document}
